\section{An equivalent crossing problem in weight blocks over every field}
\label{sec:weight-block-gate}

There is a coordinate projection that restricts the crossing search
over every field, including characteristic two. Its rank overhead is
polynomial for one weight and subexponential for a suitable collection
of weights. It uses integer counts of binary coordinates; it does not
average over the points of a finite-field group.

Let \(U=(\kk^2)^{\otimes k}\), let \(J e_a=e_{\bar a}\) complement
the binary string \(a\), and put
\[
 \Gamma_J=\operatorname{id}\otimes(X\mapsto JX^{\mathsf T}J).
\]
As before, conjugating the second summand by \(I\otimes J\) preserves
its rank and changes between ordinary and twisted partial transpose.
Thus the unrestricted minimum for \(I=M+\Gamma_J(N)\) is
\(P_\kk(2^k)\) from \eqref{eq:pt-crossing-definition}.

Partition the \(k\) coordinates into nonempty blocks \(B_i\) of sizes
\(k_i\). For \(h=(h_i)\), \(0\leq h_i\leq2k_i\), set
\[
 W_h=\operatorname{span}\{
 e_a\otimes e_b:
 \operatorname{wt}(a|_{B_i})+\operatorname{wt}(b|_{B_i})=h_i
 \text{ for every }i\}.
\]
All weights in this definition are integers, even when the field has
positive characteristic. Let \(Q_h\) be the coordinate projection onto
\(W_h\), and define
\[
 \mathcal A_{\boldsymbol k}
   =\bigoplus_h\operatorname{End}(W_h),\qquad
 P_{\rm wt}(\boldsymbol k)
   =\min_{\substack{I=M+\Gamma_J(N)\\
                    M,N\in\mathcal A_{\boldsymbol k}}}
      \bigl(\operatorname{rank}M+\operatorname{rank}N\bigr).
\]

\begin{proposition}[Weight-block projection]
\label{prop:weight-block-gate}
Over every field,
\begin{equation}\label{eq:weight-block-cost}
 P_\kk(2^k)\leq P_{\rm wt}(\boldsymbol k)
 \leq\left(\prod_i(2k_i+1)\right)P_\kk(2^k).
\end{equation}
The admissible algebra has dimension
\begin{equation}\label{eq:weight-block-dimension}
 \dim\mathcal A_{\boldsymbol k}
   =\prod_i\binom{4k_i}{2k_i}<2^{4k}.
\end{equation}
For one block the rank overhead is \(2k+1\).
\end{proposition}
\begin{proof}
Put \(\mathcal E(X)=\sum_hQ_hXQ_h\). This projection fixes \(I\).
For rows \((a,b)\) and columns \((c,d)\), twisted partial transpose
takes the entry of \(X\) at
\((a,\bar d;c,\bar b)\). In a block of size \(k_i\), the condition
\[
 \operatorname{wt}(a)+\operatorname{wt}(b)
 =\operatorname{wt}(c)+\operatorname{wt}(d)
\]
is equivalent to
\[
 \operatorname{wt}(a)+k_i-\operatorname{wt}(d)
 =\operatorname{wt}(c)+k_i-\operatorname{wt}(b).
\]
The two coordinate masks therefore agree:
\(\mathcal E\Gamma_J=\Gamma_J\mathcal E\).
Applying \(\mathcal E\) to any feasible identity gives another
feasible identity whose two summands lie in
\(\mathcal A_{\boldsymbol k}\).

There are \(H=\prod_i(2k_i+1)\) nonempty weight spaces. Since the
compressed blocks have disjoint row and column spaces,
\[
 \operatorname{rank}\mathcal E(X)
 =\sum_h\operatorname{rank}(Q_hXQ_h)
 \leq H\operatorname{rank}X.
\]
This proves the upper bound in \eqref{eq:weight-block-cost};
the lower bound follows by inclusion of feasible sets.
No division or averaging is used.

In block \(i\), the weight-\(h_i\) space has dimension
\(\binom{2k_i}{h_i}\). Thus
\[
 \dim\mathcal A_{\boldsymbol k}
 =\prod_i\sum_{h_i=0}^{2k_i}\binom{2k_i}{h_i}^{\,2}
 =\prod_i\binom{4k_i}{2k_i},
\]
by Vandermonde's identity. Each factor is strictly smaller than
\(2^{4k_i}\), proving properness.
\end{proof}

\begin{example}[One coordinate]\label{ex:weight-block-one-site}
In the ordered basis \(00,01,10,11\), the three weight spaces have
dimensions \(1,2,1\). The projection keeps precisely matrices of
the form
\[
 X=\begin{pmatrix}
 a&0&0&0\\
 0&b&c&0\\
 0&d&e&0\\
 0&0&0&f
 \end{pmatrix},\qquad
 \Gamma_J(X)=\begin{pmatrix}
 b&0&0&0\\
 0&a&c&0\\
 0&d&f&0\\
 0&0&0&e
 \end{pmatrix}.
\]
The algebra dimension is six, and the general rank overhead is three.
In characteristic two, \(p=I+F\), where \(F\) interchanges the two
spinors, has only the middle \(2\)-by-\(2\) all-one block.
It has rank one and satisfies
\(p+\Gamma_J(p)=I_4\).
This is a finite identity within the restricted space; tensoring it
does not by itself retain the two allowed cuts.
\end{example}

\begin{corollary}[The same asymptotic exponent in a proper subspace]
\label{cor:weight-block-square-gate}
Choose a partition \(\boldsymbol k=\boldsymbol k(k)\) for each \(k\)
such that
\[
 \log_2\!\left(\prod_i(2k_i+1)\right)=o(k).
\]
Then
\[
 \begin{split}
 \liminf_{k\to\infty}\frac{\log_2P_{\rm wt}(\boldsymbol k(k))}{k}
 &=\liminf_{q\to\infty}\frac{\log P_\kk(q)}{\log q},\\
 \limsup_{k\to\infty}\frac{\log_2P_{\rm wt}(\boldsymbol k(k))}{k}
 &=\limsup_{q\to\infty}\frac{\log P_\kk(q)}{\log q}.
 \end{split}
\]
One block satisfies the hypothesis. So does the partition with all
blocks of size \(\lceil\sqrt k\rceil\), except possibly its final block.
Assuming the price-space input of
Theorem~\ref{thm:square-pt-exponent}, the square is a price precisely
when the restricted exponent has limit \(\sqrt2\).
If the square is not a price, there is a \(\delta>0\) such that
\[
 P_{\rm wt}(\boldsymbol k(k))
 \leq2^{(\sqrt2-\delta)k}
 \quad\text{for all sufficiently large }k.
\]
\end{corollary}
\begin{proof}
Taking logarithms in \eqref{eq:weight-block-cost} shows that the two
ratios at \(q=2^k\) differ by \(o(1)\); this proves equality of their
lower and upper limits separately.
Restriction of all four index sets shows that \(P_\kk(q)\) is
nondecreasing. For \(2^k\leq q<2^{k+1}\),
\[
 P_\kk(2^k)\leq P_\kk(q)\leq P_\kk(2^{k+1}).
\]
The bounds in Lemma~\ref{lem:pt-upper} and \(k/(k+1)\to1\) transfer
both limits from powers of two to all integer \(q\).

For \(b=\lceil\sqrt k\rceil\), the rank overhead satisfies
\[
 \log_2 H\leq
 \left\lceil\frac{k}{b}\right\rceil\log_2(2b+1)
 =O(\sqrt k\log k)=o(k).
\]
Finally apply Theorem~\ref{thm:square-pt-exponent}. Any fixed eventual
saving there absorbs the \(o(k)\) overhead, with, for example, half
the original saving. Conversely a restricted saving is already an
unrestricted saving. This argument does not assume
submultiplicativity of \(P_\kk(q)\).
\end{proof}

Combining these limit equalities with
Theorem~\ref{thm:crossing-exponent-dichotomy} gives the exact restricted
exponent one on the nonprice branch.

This is a smaller admissible operator space at the same wire dimension,
including in characteristic two. For growing block sizes its dimension
still has the same leading exponential rate as the unrestricted space.
It supplies neither an improving decomposition nor a decision about the
square. For one block over \(\mathbb C\), the full
\(\mathrm{SL}_2\) centralizer restriction of
Section~\ref{sec:centralizer-gates} is stronger. No full
\(\mathrm{SL}_2\)-equivariance assertion in characteristic two is made.
